\documentclass[10pt,a4paper]{amsart}
\usepackage[margin=19mm]{geometry}
\usepackage{amsmath,amssymb,mathtools}
\usepackage[scr=rsfs,cal=euler]{mathalfa}
\usepackage{xfrac}
\usepackage[unicode,hidelinks,bookmarks=false]{hyperref}
\newcommand{\ie}{i.e.\ }
\newcommand{\cf}{cf.\ }
\newcommand{\resp}{respectively\ }
\newcommand{\Subsection}{\subsection}
\providecommand{\nobreakdash}{\nobreak\hbox{-}}
\setlength{\emergencystretch}{2em}
\multlinegap0pt
\usepackage{mathtools}
\usepackage{booktabs}
\usepackage{array}
\usepackage{longtable}
\usepackage{multirow}
\usepackage{enumitem}
\usepackage{setspace}
\hypersetup{
	colorlinks=true,
	linkcolor=blue,
	citecolor=blue,
	urlcolor=blue
}
\newtheorem{definition}{Definition}[section]
\newtheorem{theorem}[definition]{Theorem}
\newtheorem{lemma}[definition]{Lemma}
\newtheorem{proposition}[definition]{Proposition}
\newtheorem{remark}[definition]{Remark}
\newtheorem{corollary}[definition]{Corollary}
\begin{document}
\title[Mechanical Implications of the Teleportation of Rigid Extend]{Mechanical Implications of the Teleportation of Rigid Extended Bodies from the Perspective of Classical Mechanics and Mathematics}
\author{Robledo Maks Miranda Sette}
\date{}
\maketitle
\noindent\emph{Source and licence.} Mechanical Implications of the Teleportation of Rigid Extended Bodies from the Perspective of Classical Mechanics and Mathematics. Version arXiv:2608.03827v1 (CC BY 4.0). This material is distributed under the Creative Commons Attribution 4.0 International licence, http://creativecommons.org/licenses/by/4.0/, as recorded in the arXiv OAI licence metadata for this exact version and shown on the abstract page https://arxiv.org/abs/2608.03827v1. Full rights notice: licenses/arxiv-2608.03827-CC-BY-4.0.txt.\par\smallskip
\noindent\textit{Editorial scope.} Counted unit: the original \texttt{theorem} environment \texttt{thm:impulse} at source lines 393--410 together with its complete proof at lines 412--432; the delivered document keeps source lines 294--433 so that every referenced label and every citation resolves inside it. Native editable source is the paper's own concatenated TeX; the AMS wrapper is editorial formatting only. No publisher class, upstream script or unrelated artwork is redistributed.
	\section{A Mathematical Model for Teleportation}
	\label{sec:model}
	
	The starting point is the trajectory of a material point $P$ of mass
	$m$. Away from any teleportation event the motion is smooth, so $P$
	follows a curve of class $C^2$. We now make precise what it means for
	such a trajectory to contain a teleportation event. The object
	carrying a possible jump is taken as primitive, and the ordinary
	continuous trajectory is recovered as the special case in which no jump
	occurs.
	
	\begin{definition}[Trajectory with a teleportation instant]
		\label{def:traj}
		Let $I\subset\mathbb R$ be a nonempty open interval and let
		$t_0\in I$. A \emph{trajectory with teleportation instant $t_0$} is
		a map
		\[
		\gamma_P:I\setminus\{t_0\}\longrightarrow\mathbb R^3
		\]
		of class $C^2$ on $I\setminus\{t_0\}$ such that the one-sided limits
		\[
		P_{-}:=\lim_{t\to t_0^{-}}\gamma_P(t),
		\qquad
		P_{+}:=\lim_{t\to t_0^{+}}\gamma_P(t)
		\]
		exist and are finite, and such that the one-sided velocity limits
		\[
		\mathbf v_{-}:=\lim_{t\to t_0^{-}}\dot\gamma_P(t),
		\qquad
		\mathbf v_{+}:=\lim_{t\to t_0^{+}}\dot\gamma_P(t)
		\]
		exist and are finite.
	\end{definition}
	
	The instant $t_0$ is deliberately excluded from the domain: the point
	$P$ has no position at $t_0$.
	
	\begin{definition}[Teleportation]
		\label{def:teleport}
		The particle $P$ \emph{teleports} at $t_0$ if $P_{+}\neq P_{-}$. The
		vector $\boldsymbol\Delta:=P_{+}-P_{-}\in\mathbb R^3\setminus
		\{\mathbf 0\}$ is the \emph{teleportation displacement}, and
		$P_{-}$, $P_{+}$ are the \emph{departure} and \emph{arrival}
		points.
	\end{definition}
	
	\begin{remark}[Recovery of ordinary motion]
		If $P_{+}=P_{-}$ and $\mathbf v_{+}=\mathbf v_{-}$, then $\gamma_P$
		extends to a $C^{1}$ curve on all of $I$ by setting
		$\gamma_P(t_0):=P_{-}$; no teleportation occurs and ordinary
		continuous motion is recovered. Thus Definition~\ref{def:traj}
		contains ordinary motion as the special case
		$\boldsymbol\Delta=\mathbf 0$, $\Delta\mathbf v=\mathbf 0$: the
		trajectory with a possible jump is primitive, and continuity is a
		property that may or may not hold. On its domain $I\setminus\{t_0\}$
		the map $\gamma_P$ is everywhere differentiable; the question of
		differentiability \emph{at} $t_0$ does not arise, since
		$t_0\notin\operatorname{dom}\gamma_P$.
	\end{remark}
	
	\begin{remark}[The event is exogenous to the dynamics]
		\label{rem:exogenous}
		Throughout, Newton's law is imposed only on $I\setminus\{t_0\}$,
		and the event at $t_0$ is treated as an operation external to
		Newtonian dynamics. We do \emph{not} embed $\gamma_P$ as a
		distribution on all of $I$ and differentiate through $t_0$. This is
		deliberate. Were the discontinuous trajectory inserted into
		$m\,D^{2}\mathbf x=\mathbf F$ in the sense of distributions, with a
		position jump $\boldsymbol\Delta$ and a velocity jump
		$\Delta\mathbf v$ at $t_0$, one would obtain
		\[
		m\,D^{2}\mathbf x
		=m\,\mathbf a_{\mathrm{reg}}
		+m\,\Delta\mathbf v\,\delta_{t_0}
		+m\,\boldsymbol\Delta\,\delta'_{t_0},
		\]
		where $\delta_{t_0}$ is the Dirac distribution at $t_0$ and
		$\delta'_{t_0}$ its derivative. The velocity jump contributes an
		ordinary impulse $m\,\Delta\mathbf v\,\delta_{t_0}$; the position
		jump contributes the strictly more singular term
		$m\,\boldsymbol\Delta\,\delta'_{t_0}$, which nevertheless delivers
		\emph{zero} net linear impulse, since $\int_{\mathbb R}\delta'_{t_0}
		=0$. Removing $t_0$ from the domain does not by itself dispose of
		this term, because the associated distribution depends only on the
		$L^{1}_{\mathrm{loc}}$ class of $\gamma_P$. We therefore do not
		claim that the position jump is dynamically neutral within a
		distributional reading; we simply place the event outside the
		dynamics and compare the one-sided states, reading off the net
		impulse the mechanism must supply as the momentum jump
		$\mathbf p_{+}-\mathbf p_{-}$.
	\end{remark}
	
	It is essential to distinguish the two discontinuities. The
	\emph{position} jump $\boldsymbol\Delta=P_{+}-P_{-}$ is the
	teleportation itself; the \emph{velocity} jump
	$\Delta\mathbf v:=\mathbf v_{+}-\mathbf v_{-}$ may or may not vanish,
	and, as we now show, it is the velocity jump that fixes the net
	impulse.
	

	\begin{theorem}[Net impulse of the event]
		\label{thm:impulse}
		Let $P$ be a particle of constant mass $m>0$ following a trajectory
		with teleportation instant $t_0$ in the sense of
		Definition~\ref{def:traj}. Extend the momentum
		$\mathbf p=m\,\dot\gamma_P$ to a function of bounded variation on
		$I$ by assigning it its one-sided limits $\mathbf p_{\pm}=m\,
		\mathbf v_{\pm}$, and let $\mathbf F:=D\mathbf p$ be its
		distributional derivative, a finite vector-valued measure on
		compact subintervals. Then the atom of $\mathbf F$ at $t_0$ is
		\[
		\mathbf F(\{t_0\})=\mathbf p_{+}-\mathbf p_{-}=m\,\Delta\mathbf v,
		\]
		so the net linear impulse concentrated at the event is
		$\mathbf J=m\,\Delta\mathbf v$. In particular the event transfers a
		nonzero net linear impulse if and only if $\Delta\mathbf v\neq
		\mathbf 0$.
	\end{theorem}

	\begin{proof}
		By hypothesis the one-sided velocity limits exist and are finite, so
		$\mathbf p$ is bounded and of bounded variation on any compact
		subinterval of $I$ containing $t_0$: it is $C^{1}$ on each side of
		$t_0$ and has a single jump there. The distributional derivative of
		a $BV$ function decomposes into an absolutely continuous part, a
		jump (atomic) part and a Cantor part; here the Cantor part vanishes
		and the only atom sits at $t_0$, with mass equal to the jump
		$\mathbf p_{+}-\mathbf p_{-}=m(\mathbf v_{+}-\mathbf v_{-})
		=m\,\Delta\mathbf v$, which is nonzero precisely when
		$\Delta\mathbf v\neq\mathbf 0$ because $m>0$. Equivalently, for the
		regularized transitions $\mathbf F_\varepsilon$ obtained by
		smoothing $\mathbf p$ across $(t_0-\varepsilon,t_0+\varepsilon)$,
		\[
		\int_{t_0-\varepsilon}^{t_0+\varepsilon}\mathbf F_\varepsilon(t)\,dt
		=\mathbf p(t_0+\varepsilon)-\mathbf p(t_0-\varepsilon)
		\xrightarrow[\varepsilon\to0^{+}]{}
		\mathbf p_{+}-\mathbf p_{-}=m\,\Delta\mathbf v,
		\]
		so the impulse concentrated at $t_0$ is $m\,\Delta\mathbf v$.
	\end{proof}
	

\end{document}
